\documentclass[12pt]{article}
\usepackage[margin=1in]{geometry}
\usepackage{amsmath,amssymb,amsfonts,amsthm}
\usepackage{booktabs}
\usepackage{graphicx}
\usepackage{hyperref}
\usepackage{setspace}
\usepackage{caption}
\usepackage{subcaption}
\usepackage{tabularx}
\usepackage{enumitem}
\usepackage{microtype}
\usepackage{natbib}

\hypersetup{
    colorlinks=true,
    linkcolor=blue,
    citecolor=blue,
    urlcolor=blue
}

\onehalfspacing

\title{\textbf{Supplementary Appendix} \\ \Large Causal Emergence in U.S. Equity Industry Portfolios: Dynamic Organization and Effective Dimensionality During Systemic Stress}
\author{}
\date{}

\begin{document}
\maketitle

\tableofcontents
\newpage

\section{Scale-Invariance Numerical Test and Formal Derivation}
\label{app:scale_invariance}

A central property of the empirical estimator developed in this paper is exact invariance to common/global scalar rescaling of asset return units ($X \to c X$ for $c > 0$). This ensures that whether returns are expressed in decimals, percentages, or basis points, the information-theoretic quantities $\mathrm{CEFI}_t$ and $q_t^*$ remain numerically identical.

\subsection{Mathematical Derivation of Scale Invariance}
Consider the transition dynamics $\mathbf{x}_{t+1} = \mathbf{A}\mathbf{x}_t + \boldsymbol{\varepsilon}_{t+1}$ with $\boldsymbol{\varepsilon}_{t+1} \sim \mathcal{N}(\mathbf{0}, \boldsymbol{\Sigma}_\varepsilon)$ and unconditional state covariance $\boldsymbol{\Sigma}_x = \operatorname{Cov}(\mathbf{x}_t)$. Under a global scalar change of return units $\tilde{\mathbf{x}}_t = c \mathbf{x}_t$ ($c > 0$):
\begin{enumerate}
    \item \textbf{State Covariance:} $\tilde{\boldsymbol{\Sigma}}_x = \operatorname{Cov}(c \mathbf{x}_t) = c^2 \boldsymbol{\Sigma}_x$.
    \item \textbf{Transition Matrix Estimator:} Under the trace-scaled ridge estimator,
    \begin{equation}
    \tilde{\lambda}_t = \lambda_0 \cdot \frac{\operatorname{Tr}(\tilde{\mathbf{X}}_{\text{lag}}^\top \tilde{\mathbf{X}}_{\text{lag}})}{p} = c^2 \lambda_t
    \end{equation}
    Consequently,
    \begin{equation}
    \hat{\tilde{\mathbf{A}}} = \tilde{\mathbf{X}}_{\text{lead}}^\top \tilde{\mathbf{X}}_{\text{lag}} \left( \tilde{\mathbf{X}}_{\text{lag}}^\top \tilde{\mathbf{X}}_{\text{lag}} + \tilde{\lambda}_t \mathbf{I}_p \right)^{-1} = c^2 \mathbf{X}_{\text{lead}}^\top \mathbf{X}_{\text{lag}} \left( c^2 \mathbf{X}_{\text{lag}}^\top \mathbf{X}_{\text{lag}} + c^2 \lambda_t \mathbf{I}_p \right)^{-1} = \hat{\mathbf{A}}
    \end{equation}
    Thus, $\hat{\mathbf{A}}$ is strictly scale-invariant.
    \item \textbf{Innovation Covariance:} Residuals transform as $\tilde{\boldsymbol{\varepsilon}}_t = c \boldsymbol{\varepsilon}_t$, so under Ledoit-Wolf analytical shrinkage, $\tilde{\boldsymbol{\Sigma}}_\varepsilon = c^2 \boldsymbol{\Sigma}_\varepsilon$.
    \item \textbf{Intervention Scale:} The energy-scaled variance is:
    \begin{equation}
    \tilde{\sigma}_{do,t}^2 = \kappa^2 \cdot \frac{\operatorname{Tr}(\tilde{\boldsymbol{\Sigma}}_{x,t})}{p} = c^2 \sigma_{do,t}^2
    \end{equation}
    \item \textbf{Effective Information and Whitened SPD Representation:} In whitened coordinates, the micro-transition covariance is represented by the symmetric positive definite (SPD) matrix:
    \begin{equation}
    \mathbf{M}_{\text{micro}} = \mathbf{I}_p + \sigma_{do,t}^2 \boldsymbol{\Sigma}_\varepsilon^{-1/2} \hat{\mathbf{A}}\hat{\mathbf{A}}^\top \boldsymbol{\Sigma}_\varepsilon^{-1/2}
    \end{equation}
    By Sylvester's determinant theorem ($\det(\mathbf{I} + \mathbf{A}\mathbf{B}) = \det(\mathbf{I} + \mathbf{B}\mathbf{A})$), $\det(\mathbf{M}_{\text{micro}}) = \det(\mathbf{I}_p + \sigma_{do,t}^2 \hat{\mathbf{A}}\hat{\mathbf{A}}^\top \boldsymbol{\Sigma}_\varepsilon^{-1})$. Under global rescaling $\tilde{\mathbf{x}}_t = c \mathbf{x}_t$:
    \begin{align}
    \widetilde{EI}(\mathbf{x}) &= \frac{1}{2} \ln \det \left( \mathbf{I}_p + \tilde{\sigma}_{do,t}^2 \boldsymbol{\Sigma}_\varepsilon^{-1/2} \hat{\tilde{\mathbf{A}}}\hat{\tilde{\mathbf{A}}}^\top \tilde{\boldsymbol{\Sigma}}_\varepsilon^{-1/2} \right) \nonumber \\
    &= \frac{1}{2} \ln \det \left( \mathbf{I}_p + (c^2 \sigma_{do,t}^2) (c^2 \boldsymbol{\Sigma}_\varepsilon)^{-1/2} \hat{\mathbf{A}}\hat{\mathbf{A}}^\top (c^2 \boldsymbol{\Sigma}_\varepsilon)^{-1/2} \right) \nonumber \\
    &= \frac{1}{2} \ln \det \left( \mathbf{M}_{\text{micro}} \right) = EI(\mathbf{x})
    \end{align}
    Numerical computation is executed via the Cholesky factor of the shrunk innovation covariance $\boldsymbol{\Sigma}_\varepsilon = \mathbf{L}\mathbf{L}^\top$, setting $\mathbf{Y} = \mathbf{L}^{-1}\hat{\mathbf{A}}$ and $\mathbf{M}_{\text{micro}} = \mathbf{I}_p + \sigma_{do,t}^2 \mathbf{Y}\mathbf{Y}^\top$, guaranteeing strict positive definiteness and numerical stability.
\end{enumerate}
Identical algebraic cancellations hold for any projected macro-system $EI_q(\mathbf{W})$, guaranteeing that $\mathrm{CEFI}_t$ and $q_t^*$ (evaluated in $\mathrm{nats/DOF}$) are strictly scale-invariant.

\paragraph{Scope Delimitation.} Invariance is claimed exclusively for common scalar multipliers ($X \to c X$). Invariance is not claimed for arbitrary asset-specific diagonal scalings ($X \to X D$).

\begin{table}[htbp]
\centering
\caption{Numerical Verification of Global Return Scale Invariance Across Four Orders of Magnitude}
\label{tab:app_scale_inv}
\begin{tabular}{lccccc}
\toprule
\textbf{Scale Multiplier ($c$)} & $\mathbf{EI_{\text{micro}}}$ & $\mathbf{EI_{\text{macro}}^*}$ & $\mathbf{CEFI}$ & $\mathbf{q^*}$ & $|\Delta \mathbf{CEFI}|$ \\
\midrule
$c = 0.01$ (Basis points / 100) & 0.11381190 & 0.02644265 & -0.00022650 & 7 & $0.00$ \\
$c = 1.00$ (Standard return units) & 0.11381190 & 0.02644265 & -0.00022650 & 7 & $< 10^{-14}$ \\
$c = 100.0$ (Percentage points)   & 0.11381190 & 0.02644265 & -0.00022650 & 7 & $< 10^{-14}$ \\
$c = 10000.0$ (Basis points)      & 0.11381190 & 0.02644265 & -0.00022650 & 7 & $< 10^{-14}$ \\
\bottomrule
\end{tabular}
\end{table}

\section{VAR Transition Matrix Orientation: Mathematical Derivation and Synthetic Test}
\label{app:var_orientation}

Let $\mathbf{x}_t \in \mathbb{R}^p$ denote a column vector of asset returns at date $t$. The theoretical first-order Vector Autoregression is defined as:
\begin{equation}
\mathbf{x}_{t+1} = \mathbf{A}\mathbf{x}_t + \boldsymbol{\varepsilon}_{t+1}, \qquad \boldsymbol{\varepsilon}_{t+1} \sim \mathcal{N}(\mathbf{0}, \boldsymbol{\Sigma}_\varepsilon)
\end{equation}
where $\mathbf{A} \in \mathbb{R}^{p \times p}$ maps state $\mathbf{x}_t$ to the conditional expectation $\mathbb{E}[\mathbf{x}_{t+1} \mid \mathbf{x}_t] = \mathbf{A}\mathbf{x}_t$.

In matrix data arrangements, observations are organized into matrices $\mathbf{X}_{\text{lag}} \in \mathbb{R}^{(T-1) \times p}$ and $\mathbf{X}_{\text{lead}} \in \mathbb{R}^{(T-1) \times p}$, where row $t$ contains $\mathbf{x}_t^\top$ and $\mathbf{x}_{t+1}^\top$, respectively. Transposing the model equation yields the row-regression representation:
\begin{equation}
\mathbf{x}_{t+1}^\top = \mathbf{x}_t^\top \mathbf{A}^\top + \boldsymbol{\varepsilon}_{t+1}^\top \implies \mathbf{X}_{\text{lead}} = \mathbf{X}_{\text{lag}} \mathbf{B} + \mathbf{E}, \quad \mathbf{B} = \mathbf{A}^\top
\end{equation}
The least-squares / ridge solution for $\mathbf{B}$ is:
\begin{equation}
\hat{\mathbf{B}} = \left( \mathbf{X}_{\text{lag}}^\top \mathbf{X}_{\text{lag}} + \lambda_t \mathbf{I}_p \right)^{-1} \mathbf{X}_{\text{lag}}^\top \mathbf{X}_{\text{lead}} = \hat{\mathbf{A}}^\top
\end{equation}
Taking the transpose recovers the column transition matrix:
\begin{equation}
\hat{\mathbf{A}} = \hat{\mathbf{B}}^\top = \mathbf{X}_{\text{lead}}^\top \mathbf{X}_{\text{lag}} \left( \mathbf{X}_{\text{lag}}^\top \mathbf{X}_{\text{lag}} + \lambda_t \mathbf{I}_p \right)^{-1}
\end{equation}

\subsection{Synthetic Non-Symmetric Recovery Audit}
To verify that the implementation receives $\mathbf{A}$ and not $\mathbf{A}^\top$, we generated $T=10,000$ synthetic observations from a strongly non-symmetric transition matrix $\mathbf{A}_{\text{true}}$ ($\|\mathbf{A}_{\text{true}} - \mathbf{A}_{\text{true}}^\top\|_F = 0.738$, spectral radius $\rho = 0.528$). Fitting via our estimator yields:
\begin{equation}
\|\hat{\mathbf{A}} - \mathbf{A}_{\text{true}}\|_F = 0.061089, \qquad \|\hat{\mathbf{A}}^\top - \mathbf{A}_{\text{true}}\|_F = 1.589664
\end{equation}
The direct estimation error is over 26 times smaller than the transposed error, confirming that the estimator recovers the correctly oriented column-model transition matrix $\mathbf{A}$, rather than $\mathbf{A}^\top$.

\section{Macro-Dynamics: Interventional Lifting and Observational Closure Diagnostic}
\label{app:macro_closure}

Let $\mathbf{W} \in \mathbb{R}^{q \times p}$ denote an orthogonal coarse-graining matrix residing on the row-Stiefel manifold $\mathcal{V}_q(\mathbb{R}^p) = \{ \mathbf{W} \in \mathbb{R}^{q \times p} : \mathbf{W}\mathbf{W}^\top = \mathbf{I}_q \}$. The macroscopic state is defined as $\mathbf{y}_t = \mathbf{W}\mathbf{x}_t$.

\subsection{Interventional Channel Construction via Canonical Lifting}
Following the causal emergence framework of \citet{liu2024exact}, an intervention on the macro-state $do(\mathbf{y}_t)$ is lifted to the micro-state via the right pseudo-inverse $do(\mathbf{x}_t) = \mathbf{W}^\dagger do(\mathbf{y}_t)$. Because $\mathbf{W}\mathbf{W}^\top = \mathbf{I}_q$, the right inverse is $\mathbf{W}^\dagger = \mathbf{W}^\top (\mathbf{W}\mathbf{W}^\top)^{-1} = \mathbf{W}^\top$.

Under the interventional lifting $do(\mathbf{x}_t) = \mathbf{W}^\top \mathbf{y}_t$, the micro-state evolves as:
\begin{equation}
\mathbf{x}_{t+1} \mid do(\mathbf{y}_t) = \mathbf{A} \mathbf{W}^\top \mathbf{y}_t + \boldsymbol{\varepsilon}_{t+1}
\end{equation}
Projecting the resulting state back into the macro-subspace via $\mathbf{W}$ defines the \emph{constructed macro interventional channel}:
\begin{equation}
\mathbf{y}_{t+1} = \mathbf{W}\mathbf{x}_{t+1} = (\mathbf{W}\mathbf{A}\mathbf{W}^\top)\mathbf{y}_t + \mathbf{W}\boldsymbol{\varepsilon}_{t+1} = \mathbf{A}_M \mathbf{y}_t + \boldsymbol{\varepsilon}_{M,t+1}
\end{equation}
where $\mathbf{A}_M = \mathbf{W}\mathbf{A}\mathbf{W}^\top$ and $\boldsymbol{\Sigma}_M = \mathbf{W}\boldsymbol{\Sigma}_\varepsilon \mathbf{W}^\top$. Defining the whitened macro transition matrix:
\begin{equation}
\mathbf{M}_M(\mathbf{W}) = \mathbf{I}_q + \sigma_{do,t}^2 \boldsymbol{\Sigma}_M^{-1/2} \mathbf{A}_M \mathbf{A}_M^\top \boldsymbol{\Sigma}_M^{-1/2}
\end{equation}
By Sylvester's determinant identity, $\det(\mathbf{M}_M(\mathbf{W})) = \det(\mathbf{I}_q + \sigma_{do,t}^2 \mathbf{A}_M \mathbf{A}_M^\top \boldsymbol{\Sigma}_M^{-1})$. The continuous macro Effective Information is evaluated as:
\begin{equation}
EI_q(\mathbf{W}) = \frac{1}{2} \ln \det \left( \mathbf{M}_M(\mathbf{W}) \right) = \frac{1}{2} \ln \det \left( \mathbf{I}_q + \sigma_{do,t}^2 \mathbf{A}_M \mathbf{A}_M^\top \boldsymbol{\Sigma}_M^{-1} \right)
\end{equation}
where numerical stability is preserved via the Cholesky factor $\boldsymbol{\Sigma}_M = \mathbf{L}_M \mathbf{L}_M^\top$, setting $\mathbf{Y}_M = \mathbf{L}_M^{-1}\mathbf{A}_M$ and $\mathbf{M}_M = \mathbf{I}_q + \sigma_{do,t}^2 \mathbf{Y}_M \mathbf{Y}_M^\top$.

\subsection{Observational Projection and Closure Error Diagnostic}
Under passive observational dynamics, the true projected process is:
\begin{equation}
\mathbf{y}_{t+1} = \mathbf{W}\mathbf{x}_{t+1} = \mathbf{W}\mathbf{A}\mathbf{x}_t + \mathbf{W}\boldsymbol{\varepsilon}_{t+1} = \mathbf{W}\mathbf{A}\mathbf{W}^\top \mathbf{y}_t + \mathbf{W}\mathbf{A}(\mathbf{I}_p - \mathbf{W}^\top \mathbf{W})\mathbf{x}_t + \mathbf{W}\boldsymbol{\varepsilon}_{t+1}
\end{equation}
The middle term represents omitted micro-state dynamics orthogonal to the macro-subspace. We define the relative observational closure error diagnostic as:
\begin{equation}
r_{\text{closure},t} = \frac{\|\mathbf{W}\mathbf{A} - \mathbf{W}\mathbf{A}\mathbf{W}^\top \mathbf{W}\|_F}{\|\mathbf{W}\mathbf{A}\|_F}
\end{equation}

\begin{table}[htbp]
\centering
\caption{Observational Closure Diagnostics at Optimal Macro-Dimension $q^*$}
\label{tab:app_closure}
\begin{tabular}{lcccc}
\toprule
\textbf{Market Episode} & \textbf{Benchmark Date} & \textbf{Optimal } $\mathbf{q^*}$ & \textbf{Closure Ratio } $\mathbf{r_{\mathrm{closure}}}$ & \textbf{Observational Status} \\
\midrule
2005 Calm Benchmark   & 2005-12-30 & 1 & 0.3915 & Structurally Open \\
2008 GFC Peak         & 2008-11-20 & 1 & 0.4350 & Structurally Open \\
2020 COVID Crash      & 2020-03-23 & 2 & 0.4826 & Structurally Open \\
2000 Dot-Com Crash    & 2001-03-20 & 6 & 0.4560 & Structurally Open \\
2022 Rate Tightening  & 2022-06-15 & 1 & 0.4637 & Structurally Open \\
\bottomrule
\end{tabular}
\begin{minipage}{\textwidth}
\vspace{4pt}
\scriptsize \emph{Notes:} The non-zero closure ratios ($r_{\mathrm{closure}} \in [0.39, 0.48] > 0$) indicate that the optimized macro representations are not observationally closed; $\mathbf{W}\mathbf{A}\mathbf{W}^\top$ is therefore interpreted throughout as the transition operator of the constructed interventional macro channel rather than an autonomous observational system.
\end{minipage}
\end{table}

Because $r_{\text{closure}}$ is strictly positive ($0.39 \le r_{\text{closure}} \le 0.48$), the macro model $\mathbf{A}_M = \mathbf{W}\mathbf{A}\mathbf{W}^\top$ is formally interpreted as a \emph{constructed macro interventional transition operator} under the canonical lifting $\mathbf{W}^\top$, rather than an autonomous closed observational projection.

\section{Optimizer Budget Sensitivity Diagnostic}
\label{app:stiefel_opt}

\subsection{Riemannian Gradient and Canonical Metric Duality on Row-Stiefel}
The row-Stiefel manifold is defined as $\mathcal{V}_q(\mathbb{R}^p) = \{ \mathbf{W} \in \mathbb{R}^{q \times p} : \mathbf{W}\mathbf{W}^\top = \mathbf{I}_q \}$. The canonical metric on $T_{\mathbf{W}}\mathcal{V}_q(\mathbb{R}^p)$ is defined by \citet{edelman1998geometry} and \citet{absil2008optimization}:
\begin{equation}
g_{\mathbf{W}}^{\text{canonical}}(\mathbf{\Delta}_1, \mathbf{\Delta}_2) = \operatorname{Tr}\left( \mathbf{\Delta}_1 \left(\mathbf{I}_p - \frac{1}{2}\mathbf{W}^\top \mathbf{W}\right) \mathbf{\Delta}_2^\top \right)
\end{equation}
Under this metric, the Riemannian gradient of a differentiable scalar function $f(\mathbf{W})$ with Euclidean gradient $\mathbf{G} = \nabla_{\mathbf{W}} f(\mathbf{W})$ is:
\begin{equation}
\operatorname{grad}_{\mathcal{R}} f(\mathbf{W}) = \mathbf{G} - \mathbf{W}\mathbf{G}^\top \mathbf{W}
\end{equation}

\paragraph{Theorem (Canonical Metric Inner Product Duality).}
For any tangent vector $\mathbf{\Delta} \in T_{\mathbf{W}}\mathcal{V}_q(\mathbb{R}^p)$ (which satisfies $\mathbf{W}\mathbf{\Delta}^\top + \mathbf{\Delta}\mathbf{W}^\top = \mathbf{0}$):
\begin{equation}
g_{\mathbf{W}}^{\text{canonical}}(\operatorname{grad}_{\mathcal{R}} f(\mathbf{W}), \mathbf{\Delta}) = \operatorname{Tr}(\mathbf{G}\mathbf{\Delta}^\top) = \langle \mathbf{G}, \mathbf{\Delta} \rangle_{\text{Euclidean}} = D f(\mathbf{W})[\mathbf{\Delta}]
\end{equation}

\subsection{Optimizer Budget Sensitivity and Calibration Diagnostic}
We evaluated 25 evenly spaced historical estimation windows comparing alternative computational budgets against a high-budget reference configuration (150 iterations, 25 multistarts). The objective gap is evaluated on the true dimension-selection criterion $J(q^*) = \frac{EI_{q^*}}{q^*} - \frac{EI_p}{p}$.

\begin{table}[htbp]
\centering
\caption{Optimizer Budget Calibration Diagnostic on Selection Objective $J(q^*)$: Candidate Configurations vs. High-Budget Reference (25/150)}
\label{tab:app_optimizer_convergence}
\resizebox{\textwidth}{!}{
\begin{tabular}{lcccccccc}
\toprule
\textbf{Configuration} & \textbf{Restarts} & \textbf{Iter} & \textbf{Cost Mult.} & \textbf{Median Gap (\%)} & \textbf{Q95 Gap (\%)} & \textbf{Pearson } $\mathbf{\rho}$ & \textbf{Exact } $\mathbf{q^*}$ & $\mathbf{q^* \pm 1}$ \\
\midrule
4/35 (Calibration baseline) & 4  & 35  & 1.00$\times$  & 25.08\% & 40.74\% & 0.8452 & 52.0\% & 88.0\% \\
8/75                        & 8  & 75  & 2.90$\times$  & 5.31\%  & 16.84\% & 0.9533 & 72.0\% & 84.0\% \\
12/100 [Production]         & 12 & 100 & 5.90$\times$  & 1.82\%  & 11.70\% & 0.9732 & 92.0\% & 100.0\% \\
16/100                      & 16 & 100 & 7.84$\times$  & 1.20\%  & 8.39\%  & 0.9756 & 88.0\% & 92.0\% \\
25/150 (Reference)          & 25 & 150 & 18.02$\times$ & 0.00\%  & 0.00\%  & 1.0000 & 100.0\% & 100.0\% \\
\bottomrule
\end{tabular}
}
\end{table}

Increasing the search budget generally reduces the objective-value gap relative to the high-budget reference. At 12 multistarts and 100 iterations, the median objective gap drops to 1.82\%, with 92.0\% exact dimensional agreement with the high-budget reference and 100.0\% agreement within $\pm 1$ dimension (Pearson $\rho = 0.9732$, Spearman $\rho_S = 0.9377$). Because observed and surrogate null series are evaluated under identical optimizer budgets, Monte Carlo hypothesis tests avoid differential optimization-budget bias.

\section{VAR Stability and Spectral Radius Diagnostics}
\label{app:spectral_radius}

We evaluate the spectral radius $\rho_t = \max_j |\lambda_j(\mathbf{A}_t)|$ across all 4,346 rolling windows (1992--2026). The empirical distribution yields:
\begin{itemize}
    \item Mean $\rho_t = 0.3207$, Median $\rho_t = 0.3078$, $Q_{95} = 0.4227$, Maximum $\rho_t = 0.6777$.
    \item Zero rolling windows exhibit $\rho_t \ge 1.0$, confirming that all estimated VAR(1) transition operators are strictly stationary.
\end{itemize}

\begin{table}[htbp]
\centering
\caption{Spectral Radius of VAR(1) Transition Matrix Across Historical Stress Episodes}
\label{tab:app_spectral_radius}
\begin{tabular}{lcccc}
\toprule
\textbf{Historical Episode} & \textbf{Mean } $\mathbf{\rho_t}$ & \textbf{Median } $\mathbf{\rho_t}$ & \textbf{Max } $\mathbf{\rho_t}$ & \textbf{Stationary (\%)} \\
\midrule
2005 Calm Market Benchmark  & 0.2948 & 0.2910 & 0.3340 & 100\% \\
2008 GFC Peak               & 0.3159 & 0.3120 & 0.4120 & 100\% \\
2020 COVID Shock            & 0.4222 & 0.4180 & 0.5890 & 100\% \\
2000 Dot-Com Crash          & 0.3024 & 0.2990 & 0.3610 & 100\% \\
2022 Rate Tightening        & 0.3015 & 0.2980 & 0.3580 & 100\% \\
\bottomrule
\end{tabular}
\end{table}

\section{VAR Dynamic Specification and Lag Adequacy Diagnostics}
\label{app:var_lag_order}

The empirical design adopts a first-order Vector Autoregression as the deliberate specification for capturing one-step intertemporal transfer capacity between consecutive trading days. Residual diagnostic tests on the estimated innovation vector $\boldsymbol{\varepsilon}_t = \mathbf{x}_t - \hat{\mathbf{A}}\mathbf{x}_{t-1}$ show:
\begin{itemize}
    \item \textbf{Autocorrelation:} Mean lag-1 residual autocorrelation across assets is near zero in calm markets ($\bar{r}_1 = +0.0064$), GFC ($\bar{r}_1 = -0.0454$), and COVID ($\bar{r}_1 = +0.0097$).
    \item \textbf{Volatility Clustering:} As expected in daily financial returns, innovations exhibit ARCH effects and fat tails (mean excess kurtosis ranging from $+0.74$ in 2005 to $+6.64$ in March 2020).
\end{itemize}

\section{Full Matched Null Hierarchy Simulation Protocol and Results}
\label{app:null_protocol}

\subsection{Surrogate Generation and Estimation Protocol}
For each historical benchmark estimation window $\mathbf{X} \in \mathbb{R}^{W \times p}$ ($W=500$ trading days, $p=30$ industry portfolios), surrogate null time series ensembles are generated and evaluated under the following unified protocol:
\begin{enumerate}
    \item \textbf{$H_0^{\text{circ}}$ (Circular Shift Null)}: Each portfolio return column is independently rolled by a random integer lag $\tau_j \in [10, W-10]$ with periodic boundary wrapping, destroying contemporaneous and cross-lag correlations while preserving marginal autocorrelation.
    \item \textbf{$H_0^{\text{diag}}$ (Decoupled Diagonal VAR Null)}: Simulates uncoupled AR(1) processes $\mathbf{x}_{t+1} = \operatorname{diag}(\hat{\mathbf{A}})\mathbf{x}_t + \boldsymbol{\varepsilon}_{t+1}$ with independent Gaussian innovations $\boldsymbol{\varepsilon}_{t+1} \sim \mathcal{N}(\mathbf{0}, \operatorname{diag}(\boldsymbol{\Sigma}_\varepsilon))$. Initial state is drawn from the innovation distribution, initialized with a 50-step burn-in to reach empirical stationarity, followed by $W=500$ simulated observations.
    \item \textbf{$H_0^{\text{static}}$ (Static Correlation Mode Null)}: Synchronously permutes time steps across all assets via a uniform random permutation $\pi \in \mathcal{S}_W$, preserving the exact instantaneous covariance $\boldsymbol{\Sigma}_x$ while destroying all intertemporal causal transitions $\mathbf{x}_t \to \mathbf{x}_{t+1}$.
    \item \textbf{$H_0^{\text{diag+contemp}}$ (Cross-Lag Network Isolation Null)}: Simulates $\mathbf{x}_{t+1} = \operatorname{diag}(\hat{\mathbf{A}})\mathbf{x}_t + \boldsymbol{\varepsilon}_{t+1}$ driven by the full empirical contemporaneous innovation covariance $\boldsymbol{\Sigma}_\varepsilon$ via Cholesky factor $\mathbf{L}$ (with regularizing trace jitter $10^{-10}\operatorname{Tr}(\boldsymbol{\Sigma}_\varepsilon)/p$). Follows a 50-step burn-in followed by $W=500$ simulated steps, isolating own-lag persistence and instantaneous shock co-movement while setting cross-lag coupling to zero ($A_{ij}=0, \forall i \neq j$).
\end{enumerate}

Every surrogate realization $\mathbf{X}_{\text{null}} \in \mathbb{R}^{W \times p}$ is passed through the complete estimation pipeline: $\hat{\mathbf{A}}_{\text{null}}$ is fitted via trace-scaled ridge regression ($\lambda_0 = 10^{-4}$), $\boldsymbol{\Sigma}_{\varepsilon,\text{null}}$ via Ledoit-Wolf shrinkage, and $\sigma_{do,\text{null}}^2 = \kappa^2 \operatorname{Tr}(\boldsymbol{\Sigma}_{x,\text{null}})/p$. The Riemannian Stiefel optimization is executed for all candidate dimensions $q \in \{1, \dots, p-1\}$ under the identical optimizer budget as observed data (12 restarts / 100 iterations), ensuring strictly symmetric Monte Carlo hypothesis testing. Reproducibility is governed by a deterministic seed schedule ($\text{seed} = 42 + \text{regime\_id} \times 10^5 + \text{null\_id} \times 10^4 + b$).

\subsection{Matched Null Inference Results}

\begin{table}[htbp]
\centering
\caption{Complete Matched Null Model Inference Table ($B=9,999$ for Primary Nulls; $B=999$ for Auxiliary Nulls)}
\label{tab:app_full_nulls}
\resizebox{\textwidth}{!}{
\begin{tabular}{llccccccc}
\toprule
\textbf{Regime} & \textbf{Null Model} & $\mathbf{B}$ & $\mathbf{CEFI_{\text{obs}}}$ & $\mathbb{E}[\mathbf{CEFI_0}]$ & $\mathbf{Q_{95}}$ & $\mathbf{z_{\text{dev}}}$ & $\mathbf{p_{\text{emp}}}$ & $\mathbf{p_{\text{Holm}}}$ \\
\midrule
\textbf{Calm Period (2005)} 
 & $H_0^{\text{circ}}$ (Circular) & 999 & 0.2023 & +0.0708 & +0.2094 & +1.94 & 0.0630 & -- \\
 & $H_0^{\text{diag}}$ (VAR Diag) & 999 & 0.2023 & +0.0128 & +0.0390 & +15.05 & 0.0010 & -- \\
 & $H_0^{\text{static}}$ (Static Mode) & 9,999 & 0.2023 & +0.1500 & +0.2253 & +1.26 & 0.1160 & 0.3480 \\
 & $H_0^{\text{diag+contemp}}$ (Cross-lag Isol.) & 9,999 & 0.2023 & +0.1654 & +0.2504 & +0.80 & 0.2059 & 0.4118 \\
\midrule
\textbf{2008 GFC Peak} 
 & $H_0^{\text{circ}}$ (Circular) & 999 & 0.3656 & +0.1808 & +0.3642 & +1.56 & 0.0490 & -- \\
 & $H_0^{\text{diag}}$ (VAR Diag) & 999 & 0.3656 & +0.0356 & +0.0691 & +18.08 & 0.0010 & -- \\
 & $H_0^{\text{static}}$ (Static Mode) & 9,999 & 0.3656 & +0.2774 & +0.3714 & +1.70 & 0.0587 & 0.2860 \\
 & $H_0^{\text{diag+contemp}}$ (Cross-lag Isol.) & 9,999 & 0.3656 & +0.2878 & +0.3695 & +1.70 & 0.0572 & 0.2860 \\
\midrule
\textbf{2020 COVID Shock} 
 & $H_0^{\text{circ}}$ (Circular) & 999 & 0.3342 & +0.1466 & +0.2640 & +2.90 & 0.0040 & -- \\
 & $H_0^{\text{diag}}$ (VAR Diag) & 999 & 0.3342 & +0.0488 & +0.0927 & +11.54 & 0.0010 & -- \\
 & $H_0^{\text{static}}$ (Static Mode) & 9,999 & 0.3342 & +0.2213 & +0.3070 & +2.40 & 0.0182 & 0.1092 \\
 & $H_0^{\text{diag+contemp}}$ (Cross-lag Isol.) & 9,999 & 0.3342 & +0.3008 & +0.4139 & +0.53 & 0.2893 & 0.4118 \\
\bottomrule
\end{tabular}
}
\end{table}

\section{Multiple-Testing Adjustments and Sensitivity Analysis}
\label{app:multiple_testing}

We define the \textbf{Primary Hypothesis Family} as the six central regime-null combinations:
\begin{equation}
\mathcal{F}_{\text{primary}} = \{ (\text{Calm}, H_0^{\text{static}}), (\text{Calm}, H_0^{\text{diag+contemp}}), (\text{GFC}, H_0^{\text{static}}), (\text{GFC}, H_0^{\text{diag+contemp}}), (\text{COVID}, H_0^{\text{static}}), (\text{COVID}, H_0^{\text{diag+contemp}}) \}
\end{equation}
Under the Holm-Bonferroni step-down procedure ($m=6$), individual nominal test rejections are evaluated alongside family-wise multiplicity-adjusted $p$-values.

\section{Causal Effective Dimension ($q^*$) vs. Static Covariance Dimensionality}
\label{app:q_vs_static_rank}

To evaluate whether $q^*$ merely mirrors static covariance concentration (Effective Rank or PCA dimension), we evaluated their relationship across all 4,346 rolling windows:
\begin{itemize}
    \item Spearman correlation between $q^*$ and Effective Rank: $\rho_S = -0.0634$.
    \item Pearson correlation between $q^*$ and Effective Rank: $\rho = -0.0781$.
    \item Spearman correlation between $q^*$ and 80\% PCA variance dimension: $\rho_S = -0.0312$.
    \item Spearman correlation between $q^*$ and 90\% PCA variance dimension: $\rho_S = -0.0450$.
    \item Overall historical concentration: 98.94\% of observations satisfy $q^* \le 4$ with modal $q^* = 1$ (accounting for 63.8\% of all windows, and 87.7\% satisfying $q^* \le 2$).
    \item Cross-method dimensional concordance: analytical SVD emergence achieves 87.7\% agreement within $\pm 1$ dimension (63.2\% exact dimensional match).
\end{itemize}
These low correlations indicate that $q^*$ is only weakly monotonically associated with static covariance dimensionality metrics.

\section{Conventional Systemic Risk Benchmarks: Collinearity and Residualized CEFI}
\label{app:collinearity_and_residuals}

\subsection{Multicollinearity Diagnostics}
Variance Inflation Factors (VIF) and condition number for the multivariate regression of $\mathrm{CEFI}_t$ on conventional proxies ($RV_t, \bar{\rho}_t, ER_t, DY_t$):
\begin{itemize}
    \item $\text{VIF}(RV) = 3.20$, $\text{VIF}(\bar{\rho}) = 27.30$, $\text{VIF}(ER) = 13.77$, $\text{VIF}(DY) = 15.11$.
    \item Condition number of the normalized design matrix: $\kappa(\mathbf{X}) = 11.80$.
\end{itemize}
Because average correlation, effective rank, and connectedness exhibit substantial multicollinearity ($\text{VIF} > 13$), individual partial regression coefficients should not be interpreted as isolated economic channels. The complete linear specification accounts for 29.31\% of linear variation ($R^2 = 29.31\%$, leaving 70.69\% unexplained by this linear combination).

\subsection{Residualized CEFI Analysis}
We construct residualized $\mathrm{CEFI}$:
\begin{equation}
\mathrm{CEFI}_{\text{res},t} = \mathrm{CEFI}_t - \hat{\mathbb{E}}[\mathrm{CEFI}_t \mid RV_t, \bar{\rho}_t, ER_t, DY_t]
\end{equation}
Evaluating $\mathrm{CEFI}_{\text{res},t}$ across episodes:
\begin{itemize}
    \item \textbf{2008 GFC Peak:} Mean $\mathrm{CEFI}_{\text{res}} = -0.0025$ (Median $= -0.0274$).
    \item \textbf{2020 COVID Shock:} Mean $\mathrm{CEFI}_{\text{res}} = +0.0017$ (Median $= +0.0115$).
    \item \textbf{2000 Dot-Com Crash:} Mean $\mathrm{CEFI}_{\text{res}} = -0.0496$ (Median $= -0.0515$).
    \item \textbf{2022 Rate Tightening:} Mean $\mathrm{CEFI}_{\text{res}} = +0.0941$ (Median $= +0.0823$).
\end{itemize}
Residualized $\mathrm{CEFI}_t$ retains episode-level variation after linear adjustment for conventional proxies.

\section{Event Study Regressions: Full HAC Lag Bandwidth Sensitivity}
\label{app:hac_bandwidth}

\begin{table}[htbp]
\centering
\caption{Sensitivity of Event Study Estimates Across Extended Newey-West Lag Bandwidths ($L=20$ to $L=250$)}
\label{tab:app_hac_sensitivity}
\begin{tabular}{cccccccc}
\toprule
\textbf{HAC Lag ($L$)} & $\mathbf{\beta_{\text{Liq}}}$ & $\mathbf{t\text{-stat}}$ & $\mathbf{\beta_{\text{Val}}}$ & $\mathbf{t\text{-stat}}$ & $\mathbf{\Delta \beta}$ & \textbf{Wald } $\mathbf{t\text{-stat}}$ & \textbf{Wald } $\mathbf{p\text{-val}}$ \\
\midrule
$L = 20$  & +0.034 & +1.36 & -0.071 & -2.90 & +0.105 & +3.11 & $0.0019$ \\
$L = 40$  & +0.034 & +1.08 & -0.071 & -2.20 & +0.105 & +2.40 & $0.0165$ \\
$L = 60$  & +0.034 & +0.99 & -0.071 & -1.90 & +0.105 & +2.12 & $0.0336$ \\
$L = 120$  & +0.034 & +0.99 & -0.071 & -1.50 & +0.105 & +1.83 & $0.0678$ \\
$L = 250$  & +0.034 & +1.26 & -0.071 & -1.33 & +0.105 & +1.71 & $0.0878$ \\
\bottomrule
\end{tabular}
\end{table}

\section{Leave-One-Episode-Out Crisis Sensitivity Analysis}
\label{app:leave_one_out}

To verify whether the historical contrast between liquidity crises and valuation repricing ($\Delta \beta = \beta_{\text{Liq}} - \beta_{\text{Val}} > 0$) is driven by a single outlier episode, we re-estimated the specification excluding each historical episode sequentially, computing the exact contrast Wald test using the full HAC covariance matrix.

\begin{table}[htbp]
\centering
\caption{Leave-One-Episode-Out Event Study Robustness with Full HAC Contrast Covariance}
\label{tab:app_leave_one_out}
\begin{tabular}{lccccc}
\toprule
\textbf{Excluded Episode} & $\mathbf{\beta_{\text{Liq}}}$ & $\mathbf{\beta_{\text{Val}}}$ & $\mathbf{\Delta \beta}$ & \textbf{Exact Wald } $\mathbf{t\text{-stat}}$ & \textbf{Exact Wald } $\mathbf{p\text{-val}}$ \\
\midrule
None (Full Sample) & +0.034 & -0.071 & +0.105 & +2.40 & $0.0165$ \\
Exclude Dot-Com Crash & +0.034 & +0.096 & -0.063 & -1.36 & $0.1725$ \\
Exclude 2008 GFC Peak & +0.036 & -0.071 & +0.107 & +2.60 & $0.0092$ \\
Exclude 2020 COVID Shock & +0.033 & -0.071 & +0.105 & +2.22 & $0.0264$ \\
Exclude 2022 Rate Tightening & +0.034 & -0.136 & +0.169 & +5.24 & $1.58 \times 10^{-7}$ \\
\bottomrule
\end{tabular}
\end{table}

In single-episode exclusions omitting GFC ($\Delta \beta = +0.107, t = 2.60$), COVID ($\Delta \beta = +0.105, t = 2.22$), or 2022 tightening ($\Delta \beta = +0.169, t = 5.24$), the contrast remains strictly positive and statistically significant. However, omitting the Dot-Com crash reduces the contrast to $\Delta \beta = -0.063$ ($t = -1.36, p = 0.1725$), demonstrating that the statistical precision and magnitude of the historical regime contrast are materially influenced by the Dot-Com episode.

\section{Sensitivity to Intervention Scale Parameter \texorpdfstring{$\kappa$}{kappa}}
\label{app:kappa_sensitivity}

\begin{table}[htbp]
\centering
\caption{Sensitivity of Causal Emergence to Dimensionless Intervention Scale $\kappa \in [0.25, 4.0]$}
\label{tab:app_kappa}
\begin{tabular}{cccccc}
\toprule
$\mathbf{\kappa}$ & \textbf{Mean } $\mathbf{CEFI}$ & \textbf{Spearman } $\mathbf{\rho_S}$ vs Baseline & \textbf{Modal } $\mathbf{q^*}$ & $\mathbf{\beta_{\text{Liq}}}$ & $\mathbf{t\text{-stat}}$ \\
\midrule
0.25 & 0.0005 & 0.688 & 28 & +0.0015 & +1.27 \\
0.50 & 0.0316 & 0.839 & 3  & +0.0205 & +1.04 \\
1.00 & 0.2375 & 1.000 & 1  & +0.0336 & +1.08 \\
2.00 & 0.5146 & 0.894 & 1  & +0.0231 & +0.68 \\
4.00 & 0.7729 & 0.789 & 1  & -0.0038 & -0.12 \\
\bottomrule
\end{tabular}
\end{table}

\section{Sensitivity to Rolling Window Length \texorpdfstring{$W$}{W}}
\label{app:window_length}

\begin{table}[htbp]
\centering
\caption{Robustness Across Rolling Estimation Window Lengths ($W \in \{500, 750, 1000\}$ Trading Days)}
\label{tab:app_window}
\begin{tabular}{ccccc}
\toprule
\textbf{Window Length ($W$)} & \textbf{Mean } $\mathbf{CEFI}$ & \textbf{Spearman } $\mathbf{\rho_S}$ vs Baseline & \textbf{Optimal Phase Lag ($\ell$)} & \textbf{Lagged Corr} \\
\midrule
$W = 500$ days (Baseline) & 0.2375 & 1.000 & 0 & 1.000 \\
$W = 750$ days            & 0.1676 & 0.819 & -10 days & 0.793 \\
$W = 1000$ days           & 0.1250 & 0.690 & -140 days & 0.680 \\
\bottomrule
\end{tabular}
\end{table}

\section{Cross-Partition and Industry-Granularity Sensitivity: Fama-French 49 Industry Portfolios}
\label{app:ff49}

To evaluate whether our empirical findings are artifacts of the 30-industry partition or persist at finer sectoral granularity, we replicate the estimation pipeline on the 49 Fama-French Industry Portfolios ($p=49$) across all $q \in \{1, \dots, 48\}$. Because this test evaluates an alternative classification of the same underlying U.S. equity universe, it represents a sensitivity check on industry granularity and cross-partition stability rather than cross-universe or out-of-sample robustness:
\begin{itemize}
    \item Sample period: 1990--2026 ($T = 9,190$ trading days, 4,346 rolling windows).
    \item Mean $\mathrm{CEFI}_{FF49} = 0.3592$, Median $\mathrm{CEFI}_{FF49} = 0.3483$, Modal $q^* = 1$.
    \item 97.13\% of historical rolling windows exhibit $q^* \le 4$.
    \item At the March 2020 benchmark, $\mathrm{CEFI}_{\text{obs}} = 0.4578$; this provides an exploratory sensitivity check showing that low causal effective dimensionality and positive emergence density persist when the microscopic system is observed at finer industrial granularity.
\end{itemize}

\section{Theoretical Cross-Method Benchmarking Formulations: Liu et al. (2024) and Liu et al. (2025)}
\label{app:cross_method_formulations}

\subsection{Liu et al. (2024) Exact Bounded-Uniform Theory}
In the exact formulation of \citet{liu2024exact}, causal emergence is developed for linear stochastic systems under uniform interventions over a bounded compact domain $\mathcal{D} = [-\delta, \delta]^p$:
\begin{equation}
\mathbf{x}_{t+1} = \mathbf{A}\mathbf{x}_t + \boldsymbol{\varepsilon}_{t+1}, \qquad \boldsymbol{\varepsilon}_{t+1} \sim \mathcal{N}(\mathbf{0}, \boldsymbol{\Sigma}_\varepsilon), \quad do(\mathbf{x}_t) \sim \mathcal{U}([-\delta, \delta]^p)
\end{equation}
The continuous effective information of the microscopic system is evaluated as:
\begin{equation}
\mathcal{J}_{\text{micro}}(\mathbf{x}) = \frac{1}{2}\ln\det\left(\mathbf{A}\boldsymbol{\Sigma}_x \mathbf{A}^\top + \boldsymbol{\Sigma}_\varepsilon\right) - \frac{1}{2}\ln\det(\boldsymbol{\Sigma}_\varepsilon)
\end{equation}
where $\boldsymbol{\Sigma}_x$ is the state covariance. For an orthogonal projection $\mathbf{W} \in \mathcal{V}_q(\mathbb{R}^p)$ onto a $q$-dimensional macro-state $\mathbf{y}_t = \mathbf{W}\mathbf{x}_t$, the macro parameters are $\mathbf{A}_M = \mathbf{W}\mathbf{A}\mathbf{W}^\top$, $\boldsymbol{\Sigma}_M = \mathbf{W}\boldsymbol{\Sigma}_\varepsilon \mathbf{W}^\top$, and $\boldsymbol{\Sigma}_Y = \mathbf{W}\boldsymbol{\Sigma}_x \mathbf{W}^\top$. The macro effective information is:
\begin{equation}
\mathcal{J}_{\text{macro}}(\mathbf{y}) = \frac{1}{2}\ln\det\left(\mathbf{A}_M \boldsymbol{\Sigma}_Y \mathbf{A}_M^\top + \boldsymbol{\Sigma}_M\right) - \frac{1}{2}\ln\det(\boldsymbol{\Sigma}_M)
\end{equation}
The normalized causal emergence metric per degree of freedom is defined as:
\begin{equation}
\Delta \mathcal{J}(\mathbf{W}) = \frac{\mathcal{J}_{\text{macro}}(\mathbf{y})}{q} - \frac{\mathcal{J}_{\text{micro}}(\mathbf{x})}{p}
\end{equation}
and the optimal macro representation is obtained by maximizing $\Delta \mathcal{J}(\mathbf{W})$ over $\mathbf{W} \in \mathcal{V}_q(\mathbb{R}^p)$ and $q \in \{1, \dots, p-1\}$.

\subsection{Liu et al. (2025) Singular Value Decomposition (SVD) Formulation}
\citet{liu2025singular} propose an analytical approach to causal emergence in Gaussian iterative dynamics that circumvents iterative manifold optimization. Let $\boldsymbol{\Sigma}_\varepsilon = \mathbf{L}\mathbf{L}^\top$ denote the Cholesky factorization of the innovation covariance. The noise-whitened transition operator is defined as:
\begin{equation}
\mathbf{K} = \mathbf{L}^{-1}\mathbf{A} \in \mathbb{R}^{p \times p}
\end{equation}
Taking the Singular Value Decomposition (SVD) of $\mathbf{K}$:
\begin{equation}
\mathbf{K} = \mathbf{U} \mathbf{S} \mathbf{V}^\top, \qquad \mathbf{S} = \operatorname{diag}(s_1, s_2, \dots, s_p), \quad s_1 \ge s_2 \ge \dots \ge s_p \ge 0
\end{equation}
The analytical coarse-graining projection for dimension $q$ is given by the leading $q$ left singular vectors: $\mathbf{W}_{\text{SVD}} = \mathbf{U}_{:, 1:q}^\top \in \mathbb{R}^{q \times p}$. Under an isotropic Gaussian intervention with variance $\sigma_{do}^2$, the micro and macro Effective Information evaluate in closed form directly from the singular spectrum:
\begin{equation}
EI_{\text{micro}} = \frac{1}{2}\sum_{i=1}^p \ln\left(1 + \sigma_{do}^2 s_i^2\right), \qquad EI_{\text{macro}}(q) = \frac{1}{2}\sum_{i=1}^q \ln\left(1 + \sigma_{do}^2 s_i^2\right)
\end{equation}
The SVD causal emergence criterion per degree of freedom is:
\begin{equation}
\Delta EI_{\text{SVD}}(q) = \frac{1}{2q}\sum_{i=1}^q \ln\left(1 + \sigma_{do}^2 s_i^2\right) - \frac{1}{2p}\sum_{i=1}^p \ln\left(1 + \sigma_{do}^2 s_i^2\right)
\end{equation}
with optimal dimension $q_{\text{SVD}}^* = \operatorname{argmax}_{q} \Delta EI_{\text{SVD}}(q)$. \citet{liu2025singular} demonstrate that SVD emergence and Stiefel-based $EI$ maximization are concordant under Gaussian assumptions. In our empirical rolling analysis, the two frameworks exhibit strong positive correlation ($\rho = 0.819$, Spearman $\rho_S = 0.844$; exact dimensional agreement 63.2\%, and 87.7\% within $\pm 1$ dimension), corroborating the numerical reliability of the Stiefel optimizer.

\section{Computational Environment and Software Manifest}
\label{app:comp_env}

All estimations were performed in the following computational environment:
\begin{itemize}
    \item \textbf{High-Performance Hardware:} Multi-threaded host architecture (16 physical cores, 32 threads) and hardware-accelerated tensor execution environments.
    \item \textbf{Operating Systems:} Linux 6.8 (production server) and macOS (Darwin 24.3.0, Apple Silicon reference-parity platform).
    \item \textbf{Python \& Core Toolchain:} Python 3.11.0, PyTorch 2.3.1 (FP64 ATen Riemannian tensor engine), NumPy 2.0.2, SciPy 1.17.1, pandas 2.3.3, statsmodels 0.14.6, scikit-learn 1.8.0.
    \item \textbf{Parallel Architecture:} Concurrent multi-process execution partitioning rolling windows and surrogate null regimes across dedicated compute streams.
    \item \textbf{Reproducibility:} Deterministic host-generated orthogonal initialization seeds; CPU and GPU tensor implementations exhibited numerical agreement to at least eight decimal places under the fixed deterministic test configuration.
\end{itemize}

\begin{thebibliography}{28}
\bibitem[Absil et~al., 2008]{absil2008optimization}
Absil, P.-A., Mahony, R., Sepulchre, R., 2008.
\newblock Optimization Algorithms on Matrix Manifolds.
\newblock Princeton University Press, Princeton, NJ. \url{https://doi.org/10.1515/9781400830244}

\bibitem[Acharya et~al., 2017]{acharya2017measuring}
Acharya, V. V., Pedersen, L. H., Philippon, T., Richardson, M., 2017.
\newblock Measuring Systemic Risk.
\newblock \emph{Review of Financial Studies} 30(1), 2--47. \url{https://doi.org/10.1093/rfs/hhw088}

\bibitem[Adrian and Brunnermeier, 2016]{adrian2016covar}
Adrian, T., Brunnermeier, M. K., 2016.
\newblock CoVaR.
\newblock \emph{American Economic Review} 106(7), 1705--1741. \url{https://doi.org/10.1257/aer.20120555}

\bibitem[Ahelegbey et~al., 2022]{ahelegbey2022network}
Ahelegbey, D. F., Cerchiello, P., Scaramozzino, R., 2022.
\newblock Network based evidence of the financial impact of Covid-19 pandemic.
\newblock \emph{International Review of Financial Analysis} 81, 102101. \url{https://doi.org/10.1016/j.irfa.2022.102101}

\bibitem[Ardakani, 2024]{ardakani2024transfer}
Ardakani, O.~M., 2024.
\newblock Portfolio optimization with transfer entropy constraints.
\newblock \emph{International Review of Financial Analysis} 96, 103644. \url{https://doi.org/10.1016/j.irfa.2024.103644}

\bibitem[Barun{\'i}k and K{\v{r}}ehl{\'i}k, 2018]{barunik2018measuring}
Barun{\'i}k, J., K{\v{r}}ehl{\'i}k, T., 2018.
\newblock Measuring the frequency dynamics of financial connectedness and systemic risk.
\newblock \emph{Journal of Financial Econometrics} 16(2), 271--296. \url{https://doi.org/10.1093/jjfinec/nby001}

\bibitem[Bekiros et~al., 2017]{bekiros2017information}
Bekiros, S., Nguyen, D. K., Uddin, G. S., Sjo, B., 2017.
\newblock Information diffusion, cluster formation and entropy-based network dynamics in equity and commodity markets.
\newblock \emph{European Journal of Operational Research} 256(3), 945--961. \url{https://doi.org/10.1016/j.ejor.2016.06.052}

\bibitem[Billio et~al., 2012]{billio2012econometric}
Billio, M., Getmansky, M., Lo, A. W., Pelizzon, L., 2012.
\newblock Econometric measures of connectedness and systemic risk in the finance and insurance sectors.
\newblock \emph{Journal of Financial Economics} 104(3), 535--559. \url{https://doi.org/10.1016/j.jfineco.2011.12.010}

\bibitem[Bouri et~al., 2021]{bouri2021return}
Bouri, E., Cepni, O., Gabauer, D., Gupta, R., 2021.
\newblock Return connectedness across asset classes around the COVID-19 outbreak.
\newblock \emph{International Review of Financial Analysis} 73, 101646. \url{https://doi.org/10.1016/j.irfa.2020.101646}

\bibitem[Brownlees and Engle, 2017]{brownlees2017srisk}
Brownlees, C., Engle, R. F., 2017.
\newblock SRISK: A conditional capital shortfall measure of systemic risk.
\newblock \emph{Review of Financial Studies} 30(1), 48--79. \url{https://doi.org/10.1093/rfs/hhw060}

\bibitem[Demirer et~al., 2018]{demirer2018estimating}
Demirer, M., Diebold, F.~X., Liu, L., Yilmaz, K., 2018.
\newblock Estimating global bank network connectedness.
\newblock \emph{Journal of Applied Econometrics} 33(1), 1--15. \url{https://doi.org/10.1002/jae.2585}

\bibitem[Diebold and Yilmaz, 2012]{diebold2012better}
Diebold, F. X., Yilmaz, K., 2012.
\newblock Better to give than to receive: Predictive directional measurement of volatility spillovers.
\newblock \emph{International Journal of Forecasting} 28(1), 57--66. \url{https://doi.org/10.1016/j.ijforecast.2011.02.006}

\bibitem[Diebold and Yilmaz, 2014]{diebold2014network}
Diebold, F. X., Yilmaz, K., 2014.
\newblock On the network topology of variance decompositions: Measuring the connectedness of financial firms.
\newblock \emph{Journal of Econometrics} 182(1), 119--134. \url{https://doi.org/10.1016/j.jeconom.2014.04.012}

\bibitem[Dimpfl and Peter, 2013]{dimpfl2013using}
Dimpfl, T., Peter, F. J., 2013.
\newblock Using transfer entropy to measure dynamic information flows between markets.
\newblock \emph{Journal of International Financial Markets, Institutions and Money} 27, 273--294. \url{https://doi.org/10.1016/j.intfin.2013.09.003}

\bibitem[Edelman et~al., 1998]{edelman1998geometry}
Edelman, A., Arias, T. A., Smith, S. T., 1998.
\newblock The geometry of algorithms with orthogonality constraints.
\newblock \emph{SIAM Journal on Matrix Analysis and Applications} 20(2), 303--353. \url{https://doi.org/10.1137/S0895479895290954}

\bibitem[Hoel, 2017]{hoel2017when}
Hoel, E. P., 2017.
\newblock When the map is better than the territory.
\newblock \emph{Entropy} 19(5), 188. \url{https://doi.org/10.3390/e19050188}

\bibitem[Hoel et~al., 2013]{hoel2013quantifying}
Hoel, E. P., Albantakis, L., Tononi, G., 2013.
\newblock Quantifying causal emergence shows that macro can beat micro.
\newblock \emph{Proceedings of the National Academy of Sciences} 110(49), 19790--19795. \url{https://doi.org/10.1073/pnas.1314922110}

\bibitem[Klein and Hoel, 2020]{klein2020emergence}
Klein, B., Hoel, E., 2020.
\newblock The emergence of informative higher scales in complex networks.
\newblock \emph{Complexity} 2020, 8932526. \url{https://doi.org/10.1155/2020/8932526}

\bibitem[Kritzman et~al., 2011]{kritzman2011principal}
Kritzman, M., Li, Y., Page, S., Rigobon, R., 2011.
\newblock Principal components as a measure of systemic risk.
\newblock \emph{Journal of Portfolio Management} 37(4), 112--126. \url{https://doi.org/10.3905/jpm.2011.37.4.112}

\bibitem[Ledoit and Wolf, 2004]{ledoit2004well}
Ledoit, O., Wolf, M., 2004.
\newblock A well-conditioned estimator for large-dimensional covariance matrices.
\newblock \emph{Journal of Multivariate Analysis} 88(2), 365--411. \url{https://doi.org/10.1016/S0047-259X(03)00096-4}

\bibitem[Liu et~al., 2024]{liu2024exact}
Liu, K., Yuan, B., Zhang, J., 2024.
\newblock An exact theory of causal emergence for linear stochastic iteration systems.
\newblock \emph{Entropy} 26(8), 618. \url{https://doi.org/10.3390/e26080618}

\bibitem[Liu et~al., 2025]{liu2025singular}
Liu, K., Pan, L., Wang, Z., Yang, M., Yuan, B., Zhang, J., 2025.
\newblock Singular-value-decomposition-based causal emergence for Gaussian iterative systems.
\newblock \emph{Physical Review E} 112(5), 054225. \url{https://doi.org/10.1103/mfct-sxn5}

\bibitem[Mensi et~al., 2021]{mensi2021spillovers}
Mensi, W., Hernandez, J. A., Yoon, S.-M., Vo, X. V., Kang, S. H., 2021.
\newblock Spillovers and connectedness between major precious metals and major currency markets: The role of frequency factor.
\newblock \emph{International Review of Financial Analysis} 74, 101672. \url{https://doi.org/10.1016/j.irfa.2021.101672}

\bibitem[Newey and West, 1987]{newey1987hypothesis}
Newey, W. K., West, K. D., 1987.
\newblock A simple, positive semi-definite, heteroskedasticity and autocorrelation consistent covariance matrix.
\newblock \emph{Econometrica} 55(3), 703--708. \url{https://doi.org/10.2307/1913610}

\bibitem[Politis and Romano, 1994]{politis1994stationary}
Politis, D. N., Romano, J. P., 1994.
\newblock The stationary bootstrap.
\newblock \emph{Journal of the American Statistical Association} 89(428), 1303--1313. \url{https://doi.org/10.1080/01621459.1994.10476870}

\bibitem[Rosas et~al., 2020]{rosas2020reconciling}
Rosas, F. E., Mediano, P. A., Jensen, H. J., Seth, A. K., Barrett, A. B., Carhart-Harris, R. L., Bor, D., 2020.
\newblock Reconciling emergences: An information-theoretic approach to identify causal emergence in multivariate data.
\newblock \emph{PLoS Computational Biology} 16(12), e1008579. \url{https://doi.org/10.1371/journal.pcbi.1008579}

\bibitem[Roy and Vetterli, 2007]{roy2007effective}
Roy, O., Vetterli, M., 2007.
\newblock The effective rank: A measure of effective dimensionality.
\newblock 15th European Signal Processing Conference (EUSIPCO 2007), 606--610.

\bibitem[Schreiber, 2000]{schreiber2000measuring}
Schreiber, T., 2000.
\newblock Measuring information transfer.
\newblock \emph{Physical Review Letters} 85(2), 461--464. \url{https://doi.org/10.1103/PhysRevLett.85.461}

\bibitem[Yang et~al., 2025]{yang2025finding}
Yang, M., Wang, Z., Liu, K., Rong, Y., Yuan, B., Zhang, J., 2025.
\newblock Finding emergence in data by maximizing effective information.
\newblock \emph{National Science Review} 12(1), nwae279. \url{https://doi.org/10.1093/nsr/nwae279}

\bibitem[Yang and Hamori, 2023]{yang2023modeling}
Yang, L., Hamori, S., 2023.
\newblock Modeling the global sovereign credit network under climate change.
\newblock \emph{International Review of Financial Analysis} 87, 102618. \url{https://doi.org/10.1016/j.irfa.2023.102618}

\bibitem[Yao and Li, 2025]{yao2025interindustry}
Yao, C.-Z., Li, Y.-L., 2025.
\newblock Structural evolution of industry association networks in Chinese stock market under major event shocks: A comparative analysis of two crises based on partial Granger causal networks.
\newblock \emph{International Review of Financial Analysis} 107, 104572. \url{https://doi.org/10.1016/j.irfa.2025.104572}

\bibitem[Yousaf et~al., 2022]{yousaf2022linkages}
Yousaf, I., Nekhili, R., Gubareva, M., 2022.
\newblock Linkages between DeFi assets and conventional currencies: Evidence from the COVID-19 pandemic.
\newblock \emph{International Review of Financial Analysis} 81, 102082. \url{https://doi.org/10.1016/j.irfa.2022.102082}
\end{thebibliography}

\end{document}
